% Part: first-order-logic
% Chapter: syntax-and-semantics
% Section: Semantic Notions

\documentclass[../../../include/open-logic-section]{subfiles}

\begin{document}

\olfileid{fol}{syn}{sem}
\olsection{అర్థపర భావనలు}

\begin{explain}
మొదటిస్థాయి భాషలకు \tetoken{నిర్మాణాలను}{!!{structure}s} నిర్వచించాక,
వాక్యాలకు సంబంధించిన కొన్ని మౌలిక అర్థపర ధర్మాలనూ, వాక్యాల మధ్య
సంబంధాలను నిర్వచించవచ్చు. వీటిలో అత్యంత సరళమైనది ఒక వాక్యానికి గల
\emph{చెల్లుబాటుతనం}. ఒక వాక్యం ప్రతి \tetoken{నిర్మాణంలో}{!!{structure}}
సంతృప్తమైతే అది చెల్లుబాటు అవుతుంది. చెల్లుబాటయ్యే వాక్యాల్లోని
తార్కికేతర సంకేతాలను ఎలా అర్థనిర్దేశం చేసినా అవి సంతృప్తమవుతాయి.
అందువల్ల చెల్లుబాటయ్యే వాక్యాలను \emph{తార్కిక సత్యాలు} అని కూడా
అంటారు. అవి ఏ \tetoken{నిర్మాణంలోనైనా}{!!{structure}} సత్యం, అంటే సంతృప్తమవుతాయి. కాబట్టి
వాటి సత్యత్వం వాటిలో కనిపించే తార్కిక సంకేతాలపైనా, వాటి
\tetoken{వాక్యనిర్మాణంపైనా}{!!{structure}} మాత్రమే ఆధారపడుతుంది. అది తార్కికేతర సంకేతాలపై లేదా వాటి
అర్థనిర్దేశంపై ఆధారపడదు.
\end{explain}

\begin{defn}[చెల్లుబాటుతనం]
ఒక వాక్యం~$!A$ \emph{చెల్లుబాటు అవుతుంది}, $\Entails !A$, అనడానికి అవసరమైన, సరిపడే షరతు ఏమిటంటే,
ప్రతి \tetoken{నిర్మాణం}{!!{structure}}~$\Struct M$కు $\Sat{M}{!A}$ నెరవేరాలి.
\end{defn}

\begin{defn}[అర్థపర అనుగమనం]
వాక్యాల సమితి~$\Gamma$ నుంచి వాక్యం~$!A$ \emph{అర్థపరంగా
అనుగమిస్తుంది}, $\Gamma \Entails !A$, అప్పుడే, అప్పుడు మాత్రమే
$\Sat{M}{\Gamma}$ అయ్యే ప్రతి \tetoken{నిర్మాణం}{!!{structure}}~$\Struct M$కు
$\Sat{M}{!A}$.
\end{defn}


\begin{defn}[సంతృప్తిపరచదగినతనం]
ఏదో ఒక \tetoken{నిర్మాణం}{!!{structure}}~$\Struct M$కు $\Sat{M}{\Gamma}$ నెరవేరితే వాక్యాల
సమితి~$\Gamma$ \emph{సంతృప్తిపరచదగినది}. $\Gamma$
సంతృప్తిపరచదగినది కాకపోతే దాన్ని \emph{సంతృప్తిపరచదగనిది} అంటాం.
\end{defn}

\begin{prop}
వాక్యం~$!A$ చెల్లుబాటు కావడానికి అవసరమైన, సరిపడే షరతు ఏమిటంటే, ప్రతి
వాక్యాల సమితి~$\Gamma$కు $\Gamma \Entails !A$.
\end{prop}

\begin{proof}
మొదటి దిశను చూపడానికి, $!A$ చెల్లుబాటు అవుతుందనీ, $\Gamma$ వాక్యాల
సమితి అనీ అనుకుందాం. $\Sat{M}{\Gamma}$ అయ్యే
\tetoken{ఒక నిర్మాణం}{!!a{structure}}~$\Struct M$ను తీసుకుందాం. $!A$ చెల్లుబాటు అవుతుంది
కాబట్టి $\Sat{M}{!A}$. అందువల్ల $\Gamma \Entails !A$.

తిరుగు దిశను ప్రతివిపర్యయం ద్వారా నిరూపించడానికి, $!A$ చెల్లుబాటు కాదని
అనుకుందాం. అప్పుడు $\Sat/{M}{!A}$ అయ్యే \tetoken{ఒక నిర్మాణం}{!!a{structure}}~$\Struct M$
ఉంటుంది. $\Gamma = \{ \ltrue \}$ అయినప్పుడు, $\ltrue$ చెల్లుబాటు
అవుతుంది కాబట్టి $\Sat{M}{\Gamma}$. అందువల్ల $\Sat{M}{\Gamma}$ కానీ
$\Sat/{M}{!A}$ అయ్యే \tetoken{ఒక నిర్మాణం}{!!a{structure}}~$\Struct M$ ఉంది; కనుక
$\Gamma$ నుంచి $!A$ అనుగమించదు.
\end{proof}

\begin{prop}
\ollabel{prop:entails-unsat}
$\Gamma \Entails !A$ అవుతుంది అప్పుడే, అప్పుడు మాత్రమే
$\Gamma \cup \{\lnot !A\}$ సంతృప్తిపరచదగనిది.
\end{prop}

\begin{proof}
మొదటి దిశలో, $\Gamma \Entails !A$ అనుకుందాం. ఇందుకు విరుద్ధంగా,
$\Sat{M}{\Gamma \cup \{ \lnot !A \}}$ అయ్యే
\tetoken{ఒక నిర్మాణం}{!!a{structure}}~$\Struct M$ ఉందని అనుకుందాం. $\Sat{M}{\Gamma}$, అలాగే
$\Gamma \Entails !A$ కాబట్టి $\Sat{M}{!A}$. అంతేకాక,
$\Sat{M}{\Gamma\cup \{ \lnot !A \}}$ కాబట్టి
$\Sat{M}{\lnot !A}$; అందువల్ల $\Sat{M}{!A}$, ~$\Sat/{M}{!A}$ రెండూ
లభిస్తాయి. ఇది వైరుధ్యం. కాబట్టి అలాంటి \tetoken{నిర్మాణం}{!!{structure}}~$\Struct M$
ఉండదు; అందువల్ల $\Gamma \cup \{ \lnot !A \}$
సంతృప్తిపరచదగనిది.

తిరుగు దిశకు, $\Gamma \cup \{ \lnot !A \}$
సంతృప్తిపరచదగనిది అనుకుందాం. కాబట్టి ప్రతి
\tetoken{నిర్మాణం}{!!{structure}}~$\Struct M$కు $\Sat/{M}{\Gamma}$ లేదా $\Sat{M}{!A}$.
అందువల్ల $\Sat{M}{\Gamma}$ అయ్యే ప్రతి \tetoken{నిర్మాణం}{!!{structure}}~$\Struct M$కు
$\Sat{M}{!A}$; కనుక $\Gamma \Entails !A$.
\end{proof}

\begin{prob}
\begin{enumerate}
\item $\Gamma \Entails \bot$ అవుతుంది అప్పుడే, అప్పుడు మాత్రమే
  $\Gamma$ సంతృప్తిపరచదగనిదని చూపండి.
\item $\Gamma \cup \{!A\} \Entails \bot$ అవుతుంది అప్పుడే, అప్పుడు
  మాత్రమే $\Gamma \Entails \lnot !A$ అని చూపండి.
\item $c$ అనేది $!A$లో గానీ $\Gamma$లో గానీ కనిపించదని అనుకుందాం.
  అప్పుడు $\Gamma \Entails \lforall[x][!A]$ అవుతుంది అప్పుడే, అప్పుడు
  మాత్రమే $\Gamma \Entails \Subst{!A}{c}{x}$ అని చూపండి.
\end{enumerate}
\end{prob}

\begin{prop}
$\Gamma \subseteq \Gamma'$, అలాగే $\Gamma \Entails !A$ అయితే,
$\Gamma' \Entails !A$.
\end{prop}

\begin{proof}
$\Gamma \subseteq \Gamma'$, అలాగే $\Gamma \Entails !A$ అని
అనుకుందాం. $\Sat{M}{\Gamma'}$ అయ్యే నిర్మాణం~$\Struct M$ను
తీసుకుందాం. అప్పుడు $\Sat{M}{\Gamma}$. $\Gamma \Entails !A$ కాబట్టి
$\Sat{M}{!A}$. అందువల్ల $\Sat{M}{\Gamma'}$ అయినప్పుడల్లా
$\Sat{M}{!A}$; కనుక $\Gamma' \Entails !A$.
\end{proof}


\begin{thm}[అర్థపర నిగమన సిద్ధాంతం]
\ollabel{thm:sem-deduction}
$\Gamma \cup \{!A\} \Entails !B$ అవుతుంది అప్పుడే, అప్పుడు మాత్రమే
$\Gamma \Entails !A \lif !B$.
\end{thm}

\begin{proof}
ఒక దిశకు, $\Gamma \cup \{ !A \} \Entails !B$ అనీ,
$\Sat{M}{\Gamma}$ అయ్యే \tetoken{ఒక నిర్మాణం}{!!a{structure}}~$\Struct M$ అనీ అనుకుందాం.
$\Sat{M}{!A}$ అయితే $\Sat{M}{\Gamma \cup \{ !A \} }$ కూడా నెరవేరుతుంది.
$\Gamma \cup \{ !A \}$ నుంచి $!B$ అనుగమిస్తుంది కాబట్టి, ఆ సందర్భంలో
$\Sat{M}{!B}$ నెరవేరుతుంది. అందువల్ల $\Sat{M}{!A \lif !B}$ నెరవేరుతుంది. కనుక
$\Gamma \Entails !A \lif !B$.

తిరుగు దిశకు, $\Gamma \Entails !A \lif !B$ అనీ,
$\Sat{M}{\Gamma \cup \{ !A \}}$ అయ్యే \tetoken{ఒక నిర్మాణం}{!!a{structure}}~$\Struct M$
అనీ అనుకుందాం. అప్పుడు $\Sat{M}{\Gamma}$; కాబట్టి
$\Sat{M}{!A \lif !B}$; అలాగే $\Sat{M}{!A}$ కాబట్టి $\Sat{M}{!B}$.
అందువల్ల $\Sat{M}{\Gamma \cup \{ !A \} }$ అయినప్పుడల్లా
$\Sat{M}{!B}$; కనుక $\Gamma \cup \{ !A \} \Entails !B$.
\end{proof}

\begin{prop}\ollabel{prop:quant-terms}
  $\Struct{M}$ ఒక \tetoken{నిర్మాణం}{!!a{structure}}, ~$!A(x)$ ఒక \tetoken{సూత్రం}{!!a{formula}},
  ~$t$ ఒక సంవృత పదం అనుకుందాం. ఆ సూత్రంలో స్వేచ్ఛా చరం ఒక్కటే ఉంది. అదే~$x$. అప్పుడు:
  \begin{tagenumerate}{prvEx,prvAll}
    \tagitem{prvEx}{$!A(t) \Entails \lexists[x][!A(x)]$}{}
    \tagitem{prvAll}{$\lforall[x][!A(x)] \Entails !A(t)$}{}
  \end{tagenumerate}
\end{prop}

\begin{proof}
  \begin{tagenumerate}{prvEx,prvAll}
    \tagitem{prvEx}{%
      \iftag{probEx}{సాధన.}{$\Sat{M}{!A(t)}$ అనుకుందాం.
        $s(x) = \Value{t}{M}$ అయ్యే ఒక చర నిర్దేశం~$s$ను తీసుకుందాం.
        $!A(t)$ ఒక \tetoken{వాక్యం}{!!a{sentence}} కాబట్టి $\Sat{M}{!A(t)}[s]$.
        \olref[ext]{prop:ext-formulas} ప్రకారం
        $\Sat{M}{!A(x)}[s]$. \olref[ass]{prop:sat-quant} ప్రకారం,
        $\Sat{M}{\lexists[x][!A(x)]}$.}}{}
    \tagitem{prvAll}{%
      \iftag{probAll}{సాధన.}{
        $\Sat{M}{\lforall[x][!A(x)]}$ అనుకుందాం.
        $s(x) = \Value{t}{M}$ అయ్యే ఒక చర నిర్దేశం~$s$ను తీసుకుందాం.
        \olref[ass]{prop:sat-quant} ప్రకారం $\Sat{M}{!A(x)}[s]$.
        \olref[ext]{prop:ext-formulas} ప్రకారం $\Sat{M}{!A(t)}[s]$.
        $!A(t)$ ఒక \tetoken{వాక్యం}{!!a{sentence}} కాబట్టి
        \olref[ass]{prop:sentence-sat-true} ప్రకారం
        $\Sat{M}{!A(t)}$.}}{}
  \end{tagenumerate}
\end{proof}

\begin{probtag}{probEx,probAll}
  \olref[fol][syn][sem]{prop:quant-terms} నిరూపణను పూర్తి చేయండి.
\end{probtag}

\end{document}
